\section{AI 技术方案}\label{sec:ai}

本章是文档的重点，依次说明 AI 技术的\kw{总体链路}、\kw{选型理由}、\kw{集成方式}与\kw{优化策略}。

\subsection{总体链路：从问题到可核验的回答}

研证不把大模型当作“直接给答案的黑盒”，而是把它嵌入一条\kw{生成—核验—复核}的闭环流水线（图~\ref{fig:qa-pipeline}）：模型负责理解与综合，程序负责核对与约束，另一个异构模型负责独立复核。

\begin{figure}[!htb]
\centering
\begin{tikzpicture}[x=1mm,y=1mm]
  % ================= (a) 流水线 =================
  \newcommand{\stage}[7]{% 序号 名字 x 样式 图标 标题 小字
    \rtnode[#4,minimum height=18mm]{#2}{#3,0}{18.4mm}{\faIcon{#5}\\[1pt]#6}{#7}
    \node[num]at([xshift=1.2mm,yshift=-1.2mm]#2.north west){#1};}
  \rtnode[boxN,minimum height=18mm]{q}{0,0}{14.6mm}{\faIcon{question-circle}\\[1pt]用户问题}{\textcolor{rtblueM}{+ 上下文}}
  \stage{1}{s1}{24.2}{box}{random}{路由与改写}{中译英检索式\\多篇逐篇展开}
  \stage{2}{s2}{49.3}{boxT}{search}{混合检索}{BM25 + 向量\\融合 + 重排}
  \stage{3}{s3}{74.4}{boxB}{brain}{综合生成}{原子主张\\逐字引文}
  \stage{4}{s4}{99.5}{boxA}{ruler-combined}{确定性核验}{身份 · 逐字\\数值 · 版本}
  \stage{5}{s5}{124.6}{boxV}{user-check}{独立复核}{支持/矛盾/不足\\阈值判定}
  \stage{6}{s6}{149.7}{boxT}{check-double}{发布}{逐句证据灯\\笔记 · 导出}
  \foreach \a/\b in {q/s1,s1/s2,s2/s3,s3/s4,s4/s5,s5/s6}\draw[flowB](\a)--(\b);
  % 模型 / 技术标签
  \node[chip=rtmuted]at(24.2,-12.8){Flash Lite};
  \node[chip=rtteal]at(49.3,-12.8){FTS5 · 本地 8B};
  \node[chip=rtblue]at(74.4,-12.8){Hydrogen};
  \node[chip=rtamber]at(99.5,-12.8){规则引擎};
  \node[chip=rtviolet]at(124.6,-12.8){Decision 1};
  \node[chip=rtteal]at(149.7,-12.8){SSE 时间线};
  \node[font=\tiny,text=rtmuted,anchor=east]at(-0.5,-12.8){执行者};
  % 回路：未通过 → 补检索 + 修订一次
  \draw[flow,draw=rtamber,line width=0.9pt,rounded corners=4pt](s4.north)--++(0,5.5)-|node[pos=0.25,above,font=\scriptsize,text=rtamber!85!black,inner sep=1.5pt]{未通过：补充检索 + 修订（至多一次）}(s2.north);
  % 账本条
  \fill[rtpaper,rounded corners=3pt](13,-22.6)rectangle(160.5,-17.2);
  \draw[rtline,rounded corners=3pt,line width=0.5pt](13,-22.6)rectangle(160.5,-17.2);
  \node[font=\scriptsize,text=rtmuted]at(86.7,-19.9){\textcolor{rtblueD}{\faIcon{receipt}}\ \ 配额账本贯穿全程：调用前事务预留 → 调用后结算；成功、失败、结果未知的调用都计账\qquad\textcolor{rtblueD}{\faIcon{save}}\ \ 每步写检查点};
  \node[font=\scriptsize\heavy,text=rtmuted,anchor=west]at(-8.5,19){(a) 问答流水线};
  % ================= (b) 输出示例 =================
  \begin{scope}[shift={(0,-32)}]
    \node[font=\scriptsize\heavy,text=rtmuted,anchor=west]at(-8,2){(b) 输出示例：逐句证据灯与原文定位};
    % 回答卡片
    \node[draw=rtline,line width=0.6pt,rounded corners=4pt,fill=white,minimum width=100mm,minimum height=39.5mm,inner sep=0pt,blur shadow={shadow blur steps=6,shadow xshift=0pt,shadow yshift=-0.6mm,shadow blur radius=1.2mm,shadow opacity=14}]at(42.5,-22.25){};
    \node[anchor=west,font=\footnotesize\heavy,text=rtnavy]at(-4.5,-7.2){\textcolor{rtblue}{\faIcon{comment-dots}}\ 为什么这篇论文训练时用了 80 张 GPU？};
    \draw[rtline,line width=0.4pt](-4.5,-10.8)--(89,-10.8);
    \newcommand{\claim}[5]{% y 颜色 状态 文本 出处
      \fill[#2](-2.6,#1)circle(1.25);
      \node[anchor=west,font=\scriptsize,text=rtink,text width=66mm,align=left,inner sep=0pt]at(0.3,#1){#4};
      \node[anchor=east,chip=#2,inner xsep=2.5pt]at(89.3,#1+1.9){#3};
      \node[anchor=east,font=\tiny,text=rtmuted,inner sep=0pt]at(89.3,#1-2.2){#5};}
    \claim{-16.2}{rtrose}{与原文矛盾}{问题前提有误：原文记载为 8 块 GPU，而非 80 块。}{v7 · 第 7 页 · §5.2}
    \claim{-24.4}{rtteal}{已核对支持}{模型在一台配有 8 块 NVIDIA P100 GPU 的机器上训练。}{v7 · 第 7 页 · §5.2}
    \claim{-32.6}{rtteal}{已核对支持}{基础模型共训练 100,000 步，约 12 小时。}{v7 · 第 7 页 · §5.2}
    \node[anchor=west,font=\tiny,text=rtmuted]at(-4.5,-38.3){\faIcon{microchip}\ gemini-4-hydrogen\quad\faIcon{user-check}\ Decision 1 复核\quad\faIcon{stream}\ 6 步执行记录};
    % 迷你 PDF 页
    \begin{scope}[shift={(114.5,-22)}]
      \node[draw=rtline,line width=0.6pt,fill=white,minimum width=30mm,minimum height=38mm,inner sep=0pt,blur shadow={shadow blur steps=6,shadow xshift=0pt,shadow yshift=-0.6mm,shadow blur radius=1.2mm,shadow opacity=16}]at(0,0){};
      \node[font=\tiny\heavy,text=rtink,anchor=west,inner sep=0pt]at(-12,15){5.2\ \ Hardware and Schedule};
      \foreach \y/\w in {11.5/24,9/23,6.5/24,4/19}\fill[rtline](-12,\y)rectangle(-12+\w,\y+0.9);
      \fill[rtamber!28](-12.6,-2.9)rectangle(12.6,1.6);
      \draw[rtamber,line width=0.7pt](-12.6,-2.9)rectangle(12.6,1.6);
      \node[font=\fontsize{4.2}{5}\selectfont,text=rtink,anchor=west,inner sep=0pt,text width=24mm,align=left]at(-12,-0.6){We trained our models on one machine with 8 NVIDIA P100 GPUs.};
      \foreach \y/\w in {-6/24,-8.5/22,-11/24,-13.5/16}\fill[rtline](-12,\y)rectangle(-12+\w,\y+0.9);
      \node[font=\tiny,text=rtmuted]at(0,-16.6){7};
      \node[font=\tiny,text=rtmuted,anchor=north]at(0,-20.5){Attention Is All You Need · v7};
    \end{scope}
    \draw[flowB,line width=1pt](93.2,-24.4)--(98.3,-24.4);
    \node[font=\tiny,text=rtblueD,anchor=south,inner sep=1pt]at(95.7,-23.4){跳转};
    % 图例
    \begin{scope}[shift={(136,-9)}]
      \node[font=\scriptsize\heavy,text=rtnavy,anchor=west,inner sep=0pt]at(0,3.5){证据灯};
      \foreach \y/\c/\t in {-2/rtteal/已核对支持,-8/rtamber/待复核,-14/rtrose/与原文矛盾,-20/rtmuted/证据不足}{
        \fill[\c](1.2,\y)circle(1.25);\node[font=\scriptsize,anchor=west,inner sep=0pt]at(4,\y){\t};}
      \node[font=\tiny,text=rtmuted,anchor=north west,text width=26mm,align=left,inner sep=0pt]at(0,-24.5){低于置信阈值的主张保持“待复核”，不会被标成支持。};
    \end{scope}
  \end{scope}
\end{tikzpicture}
\caption{研证的问答全链路。(a) 六个阶段各由最合适的模型或程序承担，核验不通过时触发至多一次补检索修订；(b) 输出的每条主张带状态灯与出处，可一键跳转到 PDF 中被高亮的原句，错误前提会被原文纠正}
\label{fig:qa-pipeline}
\end{figure}

\subsection{模型选型与理由}

我们没有采用“一个大模型包办一切”的做法，而是按任务性质把模型分为六种角色（表~\ref{tab:models}）。选型遵循三条原则：\textbf{（1）让最强的模型只做最需要推理的事}，把有限的高阶模型额度用在刀刃上；\textbf{（2）复核者与生成者异构}，避免“自己审自己”；\textbf{（3）检索模型本地部署}，数据不出域且没有调用成本。项目首版不做模型训练或微调：科研问答的瓶颈在“证据是否找得到、对得上”，而不在语言能力，这部分用检索与核验解决比微调更可控。

\begin{table}[H]
\centering\small
\caption{模型角色、选型理由与集成方式}\label{tab:models}
\begin{tabularx}{\linewidth}{P{17mm}P{35mm}LP{37mm}}
\toprule
\rthead \hd{角色} & \hd{模型（按降级顺序）} & \hd{选型理由} & \hd{集成方式}\\
\midrule
深度分析 & Gemini 4 Hydrogen → 3.6 Flash → 3.5 Flash Lite & 多文档比较、综述、Future Work 核查需要长上下文与强推理；支持按 JSON Schema 约束输出，便于逐条校验 & 原生 \code{generateContent}；结构化输出 + 本地 Pydantic 二次校验\\
\rtzebra 常规回答 & 同上角色链 & 单篇事实问答对推理要求较低，额度紧张时可平滑降级，界面如实标注实际作答模型 & 角色链配置 \code{RT\_MODEL\_ANSWER}\\
轻量批处理 & 3.5 Flash Lite → 3.6 Flash & 查询改写、任务规划、元数据抽取、逐段翻译调用量大、单次简单，需要高吞吐与高日额度 & 批量请求；逐块持久化，可断点续跑\\
\rtzebra 主张复核 & Gemini Decision 1（备用：生成式复核） & 判别式模型直接给出各选项的概率与置信度，可设阈值；独立配额，不挤占生成模型；与生成模型异构 & 原生 Decisions API；一次请求批量判定全部主张\\
向量嵌入 & \textbf{本地} Nemotron-3.5-Embed-8B（对照：Qwen3-Embedding-8B；备用：gemini-embedding-001） & 同一题集与切分下实测优于对照模型（图~\ref{fig:retrieval-eval}c）；4096 维、32K 上下文，中英跨语言检索效果好 & 自建检索运行时（BF16 推理），经鉴权回环隧道供 API 调用\\
\rtzebra 重排 & \textbf{本地} Nemotron-3.5-Rerank-8B & 交叉编码精排显著提升首位命中率（Hit@1 由 54.7\% 提升到 72.0\%） & 与嵌入模型同驻 GPU，超时自动跳过\\
\bottomrule
\end{tabularx}
\end{table}

\subsection{文档智能解析}

\subsubsection{多引擎路由与统一块结构}

论文 PDF 的形态差异极大：原生文字版、整本扫描版、含大量公式表格的教材。单一解析器无法兼顾速度与质量，因此我们设计了\kw{“原生优先、按需升级”的多引擎路由}（图~\ref{fig:parsing}）：默认使用轻量的原生解析，在 2 GB 内存内完成绝大多数论文；检测到文字层缺失或用户需要表格 / 公式结构时，再把指定页交给独立的 OCR 与版面模型服务。所有引擎的输出都归一化为同一种\kw{带类型、页码与坐标的文本块}，后续检索、核验与高亮完全不感知引擎差异。

\begin{figure}[!htb]
\centering
\begin{tikzpicture}[x=1mm,y=1mm]
  % 输入
  \rtnode[boxN,minimum height=15mm]{pdf}{0,0}{15.5mm}{\faIcon{file-pdf}\\[1pt]PDF 上传}{\textcolor{rtblueM}{内容去重}}
  \node[box,diamond,aspect=1.6,inner sep=0.5pt,font=\scriptsize,align=center](dec)at(26.5,0){文字层\\[-2pt]检测};
  \draw[flowB](pdf)--(dec);
  % 引擎
  \rtnode[boxB,minimum height=11mm]{e1}{60.5,22.5}{38mm}{原生解析 · layout-v2}{字号聚类识别标题 · 空白页确认}
  \rtnode[boxT,minimum height=11mm]{e2}{60.5,7.5}{38mm}{PaddleOCR · PP-OCRv5}{扫描件文字识别 · CPU 可用}
  \rtnode[boxA,minimum height=11mm]{e3}{60.5,-7.5}{38mm}{MinerU 4.2 版面模型}{表格 → HTML · 公式 → \LaTeX}
  \rtnode[boxV,minimum height=11mm]{e4}{60.5,-22.5}{38mm}{PaddleOCR-VL 1.6}{视觉语言模型 · 疑难区域精解}
  \draw[flowB](dec.north)|-node[pos=0.72,above,font=\scriptsize,text=rtblueD,inner sep=1.2pt]{有文字层}(e1.west);
  \draw[flow](dec.east)--++(2.5,0)|-(e2.west);
  \draw[flow](dec.east)--++(2.5,0)|-(e3.west);
  \draw[flow](dec.south)|-node[pos=0.25,left,font=\scriptsize,text=rtmuted,inner sep=1.5pt,align=right]{扫描 /\\[-2pt]复杂版面}(e4.west);
  \node[draw=rtline,dashed,rounded corners=4pt,fit=(e2)(e3)(e4),inner xsep=1.6mm,inner ysep=1.8mm](gpu){};
  \node[font=\scriptsize,text=rtmuted,anchor=north,inner sep=1pt]at([yshift=-0.4mm]gpu.south){独立解析服务 · 显存预算守护 · 按需启用};
  % 归一化
  \rtnode[box,minimum height=30mm]{norm}{99,0}{18.5mm}{\faIcon{th-large}\\[1pt]统一文本块}{类型 · 页码\\归一化坐标\\原文文本}
  \foreach \e in {e1,e2,e3,e4}\draw[flow](\e.east)--++(3.4,0)|-(norm.west);
  % 结构理解
  \rtnode[boxB,minimum height=30mm]{stru}{125,0}{18.5mm}{\faIcon{sitemap}\\[1pt]结构理解}{分级目录\\参考文献切分\\表 / 图 / 公式\\文献卡片}
  \draw[flowB](norm)--(stru);
  % 快照
  \rtnode[boxT,minimum height=30mm]{snap}{150.5,0}{17.5mm}{\faIcon{camera}\\[1pt]不可变快照}{逐页检查点\\校验后\\原子发布}
  \draw[flowB](stru)--(snap);
  % 反幻觉标注
  \node[font=\scriptsize,text=rtblueD,align=center,anchor=north]at(125,-17){\faIcon{shield-alt}\ 作者 / 机构由轻量模型配对\\[-2pt]每个字符串须在原文出现};
\end{tikzpicture}
\caption{文档智能解析：多引擎路由、统一块结构与不可变解析快照}
\label{fig:parsing}
\end{figure}

\subsubsection{关键技术点}

\begin{itemize}
  \item \textbf{逐页检查点与原子发布。}解析按页提交，541 页教材中途取消后可从断点恢复；全部页面完成并通过校验后，新快照才一次性成为“活动快照”，读者永远看不到半成品。
  \item \textbf{结构恢复。}用字号聚类识别小型大写等特殊标题；用 PDF 书签重建“3.2.2 Multi-Head Attention”这样的多级目录；参考文献逐条切分并与正文引用标记建立双向链接。
  \item \textbf{抗幻觉的元数据抽取。}作者—机构—邮箱的配对交给轻量模型完成，但\kw{模型返回的每个字符串都必须能在首页原文中找到，否则整条丢弃}，从机制上杜绝编造作者或单位。
  \item \textbf{资源隔离。}重型模型运行在独立进程中，带显存预算守护（每卡至少保留 2 GiB），超预算只终止本项目进程；主服务的 2 GB 内存承诺不受影响。
\end{itemize}

\begin{table}[H]
\centering\small
\caption{各解析引擎在赛题资料上的实测表现（CPU 环境）}\label{tab:parsers}
\begin{tabularx}{\linewidth}{P{40mm}P{28mm}rrL}
\toprule
\rthead \hd{引擎} & \hd{测试输入} & \hd{耗时} & \hd{峰值内存} & \hd{产出}\\
\midrule
原生解析（layout-v2） & 541 页教材 & 88.0 s & --- & 5,315 个文本块，含 4 页空白页确认\\
\rtzebra PaddleOCR（PP-OCRv5） & 15 页扫描版论文 & 234.4 s & 1.2 GB & 311 个文本块\\
MinerU 4.2 版面模型 & 15 页扫描版论文 & 352.7 s & 2.6 GB & 138 个结构块，含 4 张表格、6 个公式\\
\rtzebra PaddleOCR-VL 1.6 & 单页裁剪区域 & 139.2 s & 2.4 GB & 6 个精解块\\
\bottomrule
\end{tabularx}
\end{table}

\subsection{高级 RAG：混合检索与索引治理}

\subsubsection{三段式检索}

检索链路分为“\kw{召回—融合—精排}”三段（图~\ref{fig:retrieval}）。第一段用关键词与语义两条通路并行召回，互补长短：BM25 擅长专有名词与数字，向量擅长跨语言与同义改写；第二段用 RRF（Reciprocal Rank Fusion）把两路结果合并为候选池，并为多篇比较问题中的每一篇论文保留最低名额，防止某一篇“霸榜”；第三段用交叉编码重排模型对候选逐一精排，最后补入相邻的表头与条件句，保证证据块自身可读。

\begin{figure}[!htb]
\centering
\begin{tikzpicture}[x=1mm,y=1mm]
  \rtnode[boxN,minimum height=14mm]{q}{0,0}{13.5mm}{\faIcon{question-circle}\\[1pt]问题}{}
  \rtnode[box,minimum height=14mm]{rw}{22,0}{17.5mm}{查询改写}{中译英 ≤ 4 条\\逐篇展开}
  \rtnode[boxA,minimum height=11.5mm]{bm}{57,9.5}{36mm}{\faIcon{font}\ 关键词通路}{FTS5 · BM25 · 中英分词}
  \rtnode[boxT,minimum height=11.5mm]{dv}{57,-9.5}{36mm}{\faIcon{vector-square}\ 语义通路}{本地 8B 嵌入 · 余弦相似度}
  \rtnode[boxB,minimum height=14mm]{rrf}{93,0}{18mm}{RRF 融合}{候选池 Top-50\\每篇保底名额}
  \rtnode[boxV,minimum height=14mm]{rr}{118,0}{18mm}{交叉重排}{Rerank 8B\\超时自动跳过}
  \rtnode[box,minimum height=14mm]{ctx}{143,0}{18mm}{上下文补全}{相邻表头 / 条件\\元数据块直通}
  \node[solidB,circle,inner sep=1.2pt,font=\scriptsize\heavy,align=center](out)at(162.5,0){Top-K\\[-3pt]证据};
  \draw[flowB](q)--(rw);
  \draw[flowB](rw.east)--++(2.6,0)|-(bm.west);\draw[flowB](rw.east)--++(2.6,0)|-(dv.west);
  \draw[flowB](bm.east)--++(2.6,0)|-(rrf.west);\draw[flowB](dv.east)--++(2.6,0)|-(rrf.west);
  \draw[flowB](rrf)--(rr);\draw[flowB](rr)--(ctx);\draw[flowB](ctx)--(out);
  % 分段标注
  \draw[decorate,decoration={brace,amplitude=4pt,mirror},rtmuted](37.8,-18.5)--(76.2,-18.5)node[midway,below=5pt,font=\scriptsize\heavy,text=rtmuted]{① 召回：两路并行，互补长短};
  \draw[decorate,decoration={brace,amplitude=4pt,mirror},rtmuted](82.8,-18.5)--(103.2,-18.5)node[midway,below=5pt,font=\scriptsize\heavy,text=rtmuted]{② 融合};
  \draw[decorate,decoration={brace,amplitude=4pt,mirror},rtmuted](107.8,-18.5)--(153.2,-18.5)node[midway,below=5pt,font=\scriptsize\heavy,text=rtmuted]{③ 精排与补全};
\end{tikzpicture}
\caption{三段式混合检索：双通路召回、RRF 融合、交叉编码精排；语义通路不可用时自动降级为关键词模式并在界面标注}
\label{fig:retrieval}
\end{figure}

\subsubsection{检索效果评测}

我们基于赛题资料包（39 篇、23,185 个检索片段）构建了 80 道检索评测题\footnote{题目由团队通读原文后编写，含中英文改写题；5 道“无答案”题不计入检索指标。该评测衡量的是检索排序质量，不等同于最终回答的准确率。}，覆盖事实问答、541 页长文档、扫描件与表格、错误前提、版本敏感等题型，在完全相同的片段与题集上对比四种检索配置（图~\ref{fig:retrieval-eval}）。

\begin{figure}[!htb]
\centering
\begin{tikzpicture}
  % ---------- (a) 核心 ----------
  \begin{axis}[rtaxis,name=A,title={(a) 核心 4 篇（45 题）},
      xlabel={$K$},ylabel={Hit@$K$（\%）},xmin=0.5,xmax=20.5,ymin=0,ymax=104,xtick={1,5,10,15,20},ytick={0,20,...,100},
      legend to name=rtleg,legend columns=4]
    \addplot[sBMtf] table[x=k,y=sparse]{data/hitk_core.dat};\addlegendentry{BM25}
    \addplot[sDense] table[x=k,y=dense]{data/hitk_core.dat};\addlegendentry{向量（Nemotron 8B）}
    \addplot[sRRF] table[x=k,y=hybrid]{data/hitk_core.dat};\addlegendentry{BM25 + 向量（RRF）}
    \addplot[sOurs] table[x=k,y=reranked]{data/hitk_core.dat};\addlegendentry{\textbf{RRF + 重排（本系统）}}
  \end{axis}
  % ---------- (b) 拓展 ----------
  \begin{axis}[rtaxis,name=B,at={(A.east)},anchor=west,xshift=13mm,title={(b) 拓展 35 篇（30 题）},
      xlabel={$K$},xmin=0.5,xmax=20.5,ymin=0,ymax=104,xtick={1,5,10,15,20},ytick={0,20,...,100}]
    \addplot[sBMtf] table[x=k,y=sparse]{data/hitk_expansion.dat};
    \addplot[sDense] table[x=k,y=dense]{data/hitk_expansion.dat};
    \addplot[sRRF] table[x=k,y=hybrid]{data/hitk_expansion.dat};
    \addplot[sOurs] table[x=k,y=reranked]{data/hitk_expansion.dat};
    \node[font=\scriptsize\heavy,text=dvA,anchor=south east]at(axis cs:20.3,100.5){100\%};
  \end{axis}
  \node[anchor=south]at($(A.north)!0.5!(B.north)+(0,5.5mm)$){\pgfplotslegendfromname{rtleg}};
  % ---------- (c) 嵌入模型对照 ----------
  \begin{axis}[rtaxis,name=C,at={(A.south west)},anchor=north west,yshift=-17mm,width=54mm,height=50mm,
      title={(c) 嵌入模型对照（核心 45 题）},xlabel={$K$},ylabel={Hit@$K$（\%）},xmin=0.5,xmax=20.5,ymin=0,ymax=104,
      xtick={1,5,10,15,20},ytick={0,20,...,100},legend style={at={(0.99,0.03)},anchor=south east,font=\tiny,row sep=-1.5pt}]
    \addplot[sBMtf] table[x=k,y=bm25]{data/embed_compare.dat};\addlegendentry{BM25}
    \addplot[sRRF] table[x=k,y=qwen]{data/embed_compare.dat};\addlegendentry{Qwen3-Emb-8B}
    \addplot[sDense] table[x=k,y=nemotron]{data/embed_compare.dat};\addlegendentry{Nemotron-Emb-8B}
  \end{axis}
  % ---------- (d) 热力图 ----------
  \begin{scope}[shift={($(C.north east)+(45mm,-5.2mm)$)},x=1mm,y=1mm]
    \node[font=\footnotesize\heavy,text=rtnavy]at(14,10.6){(d) 各题型 Hit@5（\%）};
    \foreach \x/\t in {7.5/BM25,22.5/向量,37.5/RRF,52.5/{\textbf{+ 重排}}}\node[font=\scriptsize,text=rtink]at(\x,5.2){\t};
    \input{data/heat_hit5.tex}
    \draw[rtline,line width=0.4pt](-37,-36.3)--(69,-36.3);
  \end{scope}
\end{tikzpicture}
\caption{检索评测。(a)(b) 首个标准证据进入前 $K$ 名的题目比例：完整链路在两组题上 Hit@20 均达 100\%；(c) 相同条件下本地嵌入模型选型对照；(d) 分题型的 Hit@5，重排在长文档、错误前提与跨语言提问上收益最明显}
\label{fig:retrieval-eval}
\end{figure}

从评测中可以得到三点结论，它们直接决定了系统的默认配置：

\begin{itemize}
  \item \textbf{语义检索是必要的。}纯 BM25 在全部 75 道有答案题上的 Hit@20 仅为 70.7\%，中文提问检索英文论文时尤其吃力；加入本地向量后提升到 97.3\%。
  \item \textbf{融合负责“找全”，重排负责“排准”。}RRF 把两路候选合并后，标准证据进入候选池的覆盖更稳，但关键词噪声会拉低头部名次（Hit@1 为 36.0\%）；交叉编码重排把 Hit@1 提升到 \textbf{72.0\%}、MRR 提升到 \textbf{0.81}，并使 Hit@20 达到 \textbf{100\%}。
  \item \textbf{用项目实测选择默认嵌入模型。}相同片段与题集下，Nemotron-3.5-Embed-8B 的 Hit@5 为 95.6\%，高于 Qwen3-Embedding-8B 的 88.9\%，因此被选为默认模型。
\end{itemize}

\subsubsection{索引代际：让向量索引可升级、可回滚}

向量索引一旦与模型、切分方式绑定，后续更换模型就容易出现“新旧向量混用”的隐蔽错误。我们引入\kw{索引代际（Index Generation）}机制（图~\ref{fig:generation}）：每次构建都是一个独立代际，记录完整的模型指纹（模型 ID、固定版本、维度、精度、池化方式、提示模板与分词器哈希）；片段记录其在原文块中的精确跨度，可逐段回溯；只有覆盖率达到 100\% 并通过校验的代际才会被原子切换为活动索引，构建失败或中断的代际原样保留、可断点续建，线上检索始终使用上一个完整代际。

\begin{figure}[!htb]
\centering
\begin{tikzpicture}[x=1mm,y=1mm]
  \newcommand{\gen}[6]{\rtnode[#1,minimum height=13mm]{#2}{#3,0}{25mm}{#4}{#5}\node[chip=#6,anchor=south]at([yshift=0.8mm]#2.north){#2};}
  \rtnode[box,minimum height=13mm]{staged}{0,0}{27.5mm}{\faIcon{cut}\ 切分}{块 → 片段 + 原文跨度\\超长按 token 无损拆分}
  \rtnode[boxA,minimum height=13mm]{building}{36.5,0}{27.5mm}{\faIcon{cogs}\ 构建}{分批嵌入 · 断点续建\\失败原样保留}
  \rtnode[boxB,minimum height=13mm]{ready}{73,0}{27.5mm}{\faIcon{clipboard-check}\ 校验}{覆盖率 = 100\%\\模型指纹一致}
  \rtnode[boxT,minimum height=13mm]{active}{109.5,0}{27.5mm}{\faIcon{bolt}\ 原子激活}{成为活动索引\\旧代际保留可回滚}
  \foreach \a/\b in {staged/building,building/ready,ready/active}\draw[flowB](\a)--(\b);
  \node[chip=rtmuted,anchor=south]at([yshift=0.8mm]staged.north){staged};
  \node[chip=rtamber,anchor=south]at([yshift=0.8mm]building.north){building};
  \node[chip=rtblue,anchor=south]at([yshift=0.8mm]ready.north){ready};
  \node[chip=rtteal,anchor=south]at([yshift=0.8mm]active.north){active};
  % 指纹卡
  \node[box,anchor=west,text width=27mm,align=left,font=\scriptsize,inner sep=5pt](fp)at(131.5,0){%
    {\footnotesize\heavy\textcolor{rtblueD}{\faIcon{fingerprint}}\ 模型指纹}\\[1pt]
    \textcolor{rtmuted}{model}\ \ Nemotron 8B\\[-1pt]
    \textcolor{rtmuted}{dim}\ \ \ \ \ 4096 · BF16\\[-1pt]
    \textcolor{rtmuted}{pool}\ \ \ mean · L2\\[-1pt]
    \textcolor{rtmuted}{rev}\ \ \ \ \ 781e779e…};
  \draw[flowD](fp.west)--(active.east);
  \node[font=\scriptsize,text=rtmuted]at(70,-13){实测：22,950 个原文块 → 23,185 个片段 → 23,185 条向量，覆盖率 100\%，另有 2,189 个页眉页脚等块按规则排除};
\end{tikzpicture}
\caption{索引代际生命周期：完整构建并校验后才原子切换}
\label{fig:generation}
\end{figure}

\subsection{证据核验闭环}

“附上引用”并不等于“有证据支持”。研证对每条主张执行\kw{两级核验}（图~\ref{fig:verify}）。

\textbf{第一级：确定性核验（零模型调用）。}由规则引擎逐条检查：（1）\emph{身份}——引用的文本块必须属于本次检索候选与授权版本；（2）\emph{逐字}——引文规范化后必须是该块原文的子串；（3）\emph{数值与单位}——主张中出现的数字必须能在引文中找到，科学计数法先归一化，模型名称里的数字不当作数量；（4）\emph{元数据}——作者、机构类主张必须来自文献卡片或作者所在原文块。未通过的主张连同具体原因反馈给模型，\kw{补充检索并修订一次}；仍不通过则以“证据缺口”如实呈现，绝不发布为“已支持”。

\textbf{第二级：独立语义复核。}通过第一级的主张交给判别式模型 Gemini Decision 1，在“支持 / 矛盾 / 不足”三者中给出概率与置信度。只有同时满足\kw{概率 $\ge 0.9$、置信度 $\ge 0.8$、与次高选项的概率差 $\ge 0.25$} 才采信其结论，否则保持“待复核”。

\begin{figure}[!htb]
\centering
\begin{tikzpicture}[x=1mm,y=1mm]
  % 左：流程
  \rtnode[boxB,minimum height=13mm]{c}{0,0}{19mm}{原子主张}{+ 逐字引文}
  \node[box,text width=43mm,align=left,font=\scriptsize,inner sep=5pt](det)at(40.5,0){%
    {\footnotesize\heavy\textcolor{rtamber!85!black}{\faIcon{ruler-combined}}\ 第一级 · 确定性核验}\\[1.5pt]
    \ok\ 身份：块属于候选与授权版本\\[-0.5pt]
    \ok\ 逐字：引文是原文子串\\[-0.5pt]
    \ok\ 数值单位：数字出现在引文中\\[-0.5pt]
    \ok\ 元数据：作者 / 机构有原文依据};
  \node[box,text width=37mm,align=left,font=\scriptsize,inner sep=5pt](sem)at(95,0){%
    {\footnotesize\heavy\textcolor{rtviolet}{\faIcon{user-check}}\ 第二级 · 独立语义复核}\\[1.5pt]
    判别模型逐条给出三类概率\\[-0.5pt]
    $p\ge 0.9$\ 且\ 置信度 $\ge 0.8$\\[-0.5pt]
    且\ $p_1-p_2\ge 0.25$\ 才采信\\[-0.5pt]
    \textcolor{rtmuted}{否则保持“待复核”}};
  \draw[flowB](c)--(det);\draw[flowB](det)--node[above,font=\scriptsize,text=rtmuted,inner sep=1.5pt]{通过}(sem);
  % 修订回路
  \rtbox[boxA]{fix}{40.5,-22}{45mm}{\faIcon{redo}\ 补充检索 + 修订（至多一次）}
  \draw[flow,draw=rtamber]([xshift=5mm]det.south)--node[right,font=\scriptsize,text=rtamber!85!black,inner sep=1.5pt]{未通过}([xshift=5mm]fix.north);
  \draw[flow,draw=rtamber]([xshift=-5mm]fix.north)--([xshift=-5mm]det.south);
  \rtbox[box]{gap}{0,-22}{19mm}{\textcolor{rtmuted}{\faIcon{minus-circle}}\ 证据缺口}
  \draw[flowD](fix)--(gap);
  \node[font=\scriptsize,text=rtmuted,anchor=north]at(0,-26){修订后仍未通过};
  % 右：结果
  \foreach \y/\c/\t/\d in {12/rtteal/已核对支持/采信“支持”,4/rtrose/与原文矛盾/采信“矛盾”,-4/rtamber/待复核/低于阈值,-12/rtmuted/证据不足/采信“不足”}{
    \node[box,minimum width=35mm,minimum height=6.4mm,inner sep=1pt,font=\scriptsize,anchor=west](o\y)at(126,\y){\hspace{5.5mm}\t\ \textcolor{rtmuted}{· \d}};
    \fill[\c]([xshift=3.4mm]o\y.west)circle(1.2);
    \draw[flow](sem.east)--++(4,0)|-(o\y.west);}
  % 概率条示意
  \begin{scope}[shift={(76.5,-18)}]
    \node[font=\scriptsize,text=rtmuted,anchor=west,inner sep=0pt]at(0,0){判定示意：};
    \foreach \y/\l/\w/\c in {-4/支持/31/rtteal,-8/矛盾/1.5/rtrose,-12/不足/2.5/rtmuted}{
      \node[font=\tiny,text=rtmuted,anchor=east,inner sep=0pt]at(6.5,\y){\l};
      \fill[rtline](7.5,\y-1.1)rectangle(41.5,\y+1.1);\fill[\c](7.5,\y-1.1)rectangle(7.5+\w,\y+1.1);}
    \draw[rtink,line width=0.5pt,dash pattern=on 1pt off 1pt](38.1,-14)--(38.1,-2);
    \node[font=\tiny,text=rtink,anchor=south,inner sep=1pt]at(38.1,-2){阈值 0.9};
    \node[font=\tiny,text=rtteal!80!black,anchor=west,inner sep=0pt]at(43,-4){0.94 → 采信};
  \end{scope}
\end{tikzpicture}
\caption{两级证据核验：确定性规则拦截“引而不符”，判别模型在置信阈值之上才给出结论}
\label{fig:verify}
\end{figure}

在专门设计的对抗样例上，空引文、数字被篡改（8 → 80）、绕过主张直接在正文断言三类攻击均被拦截；Decision 1 在 11 个预设标签的校准样例上判定全部符合预期。

\subsection{有界多角色智能体}

对于“帮我找这篇论文之后有没有人做过某方向”这类开放任务，系统启用智能体模式（图~\ref{fig:agent}）。它由\kw{路由器、规划器、工具执行器、综合者、核验者、复核者}六个角色协作完成，每个角色的职责与权限相互独立。与常见的自由式 Agent 不同，研证的智能体是\kw{有界的}：

\begin{itemize}
  \item \textbf{工具白名单。}模型只能调用 11 个预先注册的工具（表~\ref{tab:tools}），每个工具的参数由 Pydantic Schema 严格校验（禁止多余字段），工具在调用者的工作区权限内执行；模型没有任何宿主 Shell 能力。
  \item \textbf{步数与预算上限。}最多两轮规划、十步工具调用、一次修订；每个工具有独立的调用预算，单任务与单工作区有模型调用总上限。
  \item \textbf{观察驱动。}第二轮规划必须基于第一轮工具的真实观察结果；工具结果、文档内容在提示词中一律标注为“数据”，不能成为指令。
  \item \textbf{全程可观测。}每一步写入持久化的公开执行记录，界面以时间线实时展示“做了什么、用了什么工具、得到什么”，但不展示也不编造模型的隐藏思维链。
\end{itemize}

\begin{figure}[!htb]
\centering
\begin{tikzpicture}[x=1mm,y=1mm]
  % 环
  \rtnode[boxN,minimum height=13mm]{task}{0,0}{17mm}{\faIcon{bullseye}\\[1pt]研究任务}{}
  \rtnode[box,minimum height=13mm]{route}{26,0}{17mm}{路由器}{任务分类\\Decision 1}
  \rtnode[boxB,minimum height=13mm]{plan}{52,0}{17mm}{规划器}{≤ 2 轮\\输出工具序列}
  \rtnode[boxT,minimum height=13mm]{act}{78,0}{17mm}{工具执行器}{≤ 10 步\\Schema 校验}
  \rtnode[boxA,minimum height=13mm]{obs}{65,-21}{25mm}{观察}{真实工具结果\\子任务等待 / 恢复}
  \rtnode[boxB,minimum height=13mm]{syn}{104,0}{17mm}{综合者}{带出处的\\原子主张}
  \rtnode[boxV,minimum height=13mm]{ver}{130,0}{17mm}{核验 + 复核}{图~\ref{fig:verify} 的\\两级闭环}
  \rtnode[boxT,minimum height=13mm]{out}{156,0}{15mm}{\faIcon{flag-checkered}\\[1pt]结果}{}
  \foreach \a/\b in {task/route,route/plan,plan/act,act/syn,syn/ver,ver/out}\draw[flowB](\a)--(\b);
  \draw[flow,rounded corners=3pt](act.south)|-(obs.east);
  \draw[flow,rounded corners=3pt](obs.west)-|(plan.south);
  \node[font=\scriptsize,text=rtmuted,fill=white,inner sep=1pt]at(65,-11){观察驱动再规划};
  % 护栏
  \node[draw=rtrose!55,dashed,rounded corners=5pt,fit=(plan)(act)(obs),inner xsep=3mm,inner ysep=3.2mm](guard){};
  \node[font=\scriptsize\heavy,text=rtrose!90!black,fill=white,inner sep=1.5pt]at(guard.north){\faIcon{shield-alt}\ 有界自治：白名单 · 步数 · 预算 · 工作区权限};
  % 工具托盘
  \begin{scope}[shift={(96.5,-22.5)}]
    \fill[rtpaper,rounded corners=4pt](0,-9.5)rectangle(69,7.5);
    \draw[rtline,line width=0.5pt,rounded corners=4pt](0,-9.5)rectangle(69,7.5);
    \node[font=\scriptsize\heavy,text=rtnavy,anchor=west]at(1.5,4.4){\faIcon{toolbox}\ 11 个白名单工具};
    \node[chip=rtblue,anchor=west]at(1.5,-0.3){\faIcon{search}\ 检索论文};
    \node[chip=rtblue,anchor=west]at(17.7,-0.3){\faIcon{download}\ 获取全文};
    \node[chip=rtblue,anchor=west]at(33.9,-0.3){\faIcon{link}\ 解析引用};
    \node[chip=rtteal,anchor=west]at(50.1,-0.3){\faIcon{file-alt}\ 解析文档};
    \node[chip=rtteal,anchor=west]at(1.5,-5.6){\faIcon{quote-left}\ 取证};
    \node[chip=rtteal,anchor=west]at(12.3,-5.6){\faIcon{code-branch}\ 版本比较};
    \node[chip=rtviolet,anchor=west]at(28.5,-5.6){\faIcon{project-diagram}\ 引用图};
    \node[chip=rtamber,anchor=west]at(42.5,-5.6){\faIcon{chart-bar}\ 统计绘图};
    \node[chip=rtmuted,anchor=west]at(57.5,-5.6){保存};
  \end{scope}
  \draw[flowD](act.south east)--(96.5,-15);
\end{tikzpicture}
\caption{有界多角色智能体：规划—执行—观察循环被白名单、步数与预算三重护栏约束}
\label{fig:agent}
\end{figure}

\begin{table}[H]
\centering\small
\caption{智能体的 11 个白名单工具}\label{tab:tools}
\begin{tabularx}{\linewidth}{P{14mm}P{43mm}LP{36mm}}
\toprule
\rthead \hd{类别} & \hd{工具} & \hd{作用} & \hd{约束}\\
\midrule
发现 & \code{search\_papers}\newline\code{resolve\_reference}\newline\code{fetch\_document} & 检索 arXiv / Crossref 元数据；把参考文献条目解析为已知作品；下载并入库全文 & 单次 ≤ 10 条候选；仅经 SSRF 安全下载器\\
\rtzebra 取证 & \code{parse\_document}\newline\code{retrieve\_evidence}\newline\code{compare\_versions} & 创建解析快照；在限定版本内检索证据块；确定性段落级版本差异 & 单次 ≤ 10 个版本、≤ 20 个证据块\\
分析 & \code{build\_citation\_graph}\newline\code{audit\_research\_question}\newline\code{compute\_csv\_stats} & 由真实引用构图并计算中心性；发起研究问题核查；校验 CSV 并计算统计量 & CSV ≤ 128 KB；统计由程序完成\\
\rtzebra 产出 & \code{render\_chart}\newline\code{save\_artifact} & 用固定模板渲染图表；保存带来源链接的研究笔记 & 不执行模型生成的代码\\
\bottomrule
\end{tabularx}
\end{table}

\textbf{Future Work 核查}是智能体能力的典型应用：从旧论文的建议出发，规划检索式 → 检索后续文献 → 自动下载并独立解析全文 → 比较发表时间窗口与实验条件 → 给出“已有直接证据 / 部分覆盖 / 当前范围证据不足”三级结论，并由独立的覆盖审查步骤复核。系统只陈述“在已检索范围内”的结论，不宣称任何方向“全球首次”。

\subsection{模型调度与配额治理}

大模型额度是评审期最稀缺的资源（例如深度模型每个项目每天仅 20 次调用）。我们把“调用哪个模型、还能调用几次”做成了一个独立的治理层（图~\ref{fig:routing}）。

\begin{figure}[!htb]
\centering
\begin{tikzpicture}[x=1mm,y=1mm]
  % 角色链
  \node[font=\scriptsize\heavy,text=rtmuted,anchor=west]at(-8.5,13){(a) 角色化降级链};
  \rtnode[boxN,minimum height=14mm]{role}{0.5,0}{17.5mm}{任务角色}{\textcolor{rtblueM}{deep · answer}\\\textcolor{rtblueM}{light · verify}}
  \rtnode[boxB,minimum height=14mm]{m1}{30,0}{17mm}{Hydrogen}{Gemini 4\\首选 · 深度推理}
  \rtnode[box,minimum height=14mm]{m2}{55.5,0}{17mm}{3.6 Flash}{第一降级\\均衡}
  \rtnode[box,minimum height=14mm]{m3}{81,0}{17mm}{Flash Lite}{第二降级\\高吞吐}
  \draw[flowB](role)--(m1);
  \draw[flow,draw=rtamber,line width=1pt](m1)--(m2);
  \draw[flow,draw=rtamber,line width=1pt](m2)--(m3);
  \node[font=\tiny,text=rtamber!85!black,anchor=south,inner sep=1.2pt]at(42.75,0.2){降级};
  \node[font=\tiny,text=rtamber!85!black,anchor=south,inner sep=1.2pt]at(68.25,0.2){降级};
  \node[font=\scriptsize,text=rtink,anchor=west,inner sep=0pt]at(-8.5,-11.2){\ok\ 仅在本地额度拒绝或上游\textbf{明确报错}时才降级};
  \node[font=\scriptsize,text=rtink,anchor=west,inner sep=0pt]at(-8.5,-15.4){\ok\ 结果\textbf{未知}的请求绝不重发，避免重复扣费};
  \node[font=\scriptsize,text=rtink,anchor=west,inner sep=0pt]at(-8.5,-19.6){\ok\ 界面如实标注实际作答模型与降级情况};
  % 账本时序
  \begin{scope}[shift={(101,0)}]
    \node[font=\scriptsize\heavy,text=rtmuted,anchor=west]at(-6,13){(b) 配额账本的事务时序};
    \foreach \x/\t in {2/Worker,30/配额账本,58/模型 API}{
      \node[box,minimum width=15mm,minimum height=5.6mm,inner sep=1pt,font=\scriptsize]at(\x,7.5){\t};
      \draw[rtline,line width=0.6pt](\x,4.7)--(\x,-20);}
    \draw[flowB](2,1)--node[above,font=\tiny,inner sep=1pt]{① 预留 RPM / TPM / RPD}(30,1);
    \draw[flowB](2,-5)--node[above,font=\tiny,inner sep=1pt,pos=0.75]{② 发送请求}(58,-5);
    \draw[flowD](58,-11)--node[above,font=\tiny,inner sep=1pt,pos=0.25]{③ 响应 / 报错 / 超时}(2,-11);
    \draw[flowB](2,-17)--node[above,font=\tiny,inner sep=1pt]{④ 结算 ok / failed / unknown}(30,-17);
  \end{scope}
\end{tikzpicture}
\caption{模型调度与配额治理：(a) 按角色配置的降级链；(b) 每次调用先预留后结算，三种结果都计入账本}
\label{fig:routing}
\end{figure}

\begin{itemize}
  \item \textbf{按项目核算。}配额按“项目 × 模型”计量（RPM、输入 TPM、每日 RPD，日界按供应商的太平洋时区），同一项目的多个密钥不会被当成多份额度。
  \item \textbf{失败也计账。}失败、超时与结果未知的请求全部计入；未知请求不自动重发，避免重复扣费与重复副作用。
  \item \textbf{响应即资产。}模型响应先落盘再做后续校验，任务重启后直接复用，不重复调用。
  \item \textbf{付费兜底默认关闭}，并保留演示储备额度，保证评审期间核心功能持续可用。
\end{itemize}

\subsection{本地模型部署}

向量嵌入与重排两个 8B 模型全部\kw{本地部署}在 GPU 主机上，由自研的检索运行时托管：以 BF16 精度推理，支持多卡分配；模型目录锁定到固定版本，并校验分词器、池化与提示模板的文件哈希，生成的“模型指纹”与索引代际一一对应。API 所在的轻量云主机通过带鉴权的回环隧道访问 GPU 主机。由于文档向量已离线算好并随索引分发，线上只有“查询向量”和“重排”需要 GPU；当 GPU 主机不可达或超时（查询 15 秒、重排 45 秒），系统自动回退到独立的云端向量代际或关键词检索，并在界面上明确标注当前检索模式。OCR 与版面模型（PaddleOCR、MinerU、PaddleOCR-VL）同样为本地离线部署，全程不依赖外部服务。

\subsection{优化策略小结}

\begin{table}[H]
\centering\small
\caption{AI 相关优化策略一览}
\begin{tabularx}{\linewidth}{P{14mm}P{38mm}L}
\toprule
\rthead \hd{目标} & \hd{策略} & \hd{做法与效果}\\
\midrule
\multirow{3}{=}{质量} & 查询改写与逐篇配额 & 中文问题改写为英文检索式；多篇比较时每篇保底名额，避免单篇垄断候选\\
 & 召回—融合—精排 & Hit@1 由 BM25 的 20.0\% 提升到 72.0\%，Hit@20 由 70.7\% 提升到 100\%\\
 & 两级核验 + 一次修订 & 规则拦截“引而不符”，判别模型阈值化复核；修订次数有上限\\
\midrule
\rtzebra & 角色化模型分工 & 深度模型只做综合推理；改写、抽取、翻译交给轻量模型；复核使用独立配额\\
\rtzebra & 确定性优先 & 哈希去重、数值统计、图算法、版本差异全部由程序完成，零模型调用\\
\rtzebra \multirow{-3}{=}{成本} & 缓存与复用 & 模型响应、翻译段落、查询向量均持久化复用；新评审空间直接复制已有向量，零重算\\
\midrule
\multirow{3}{=}{性能} & 双 Worker 池 & 文档解析与交互问答互不阻塞；541 页教材解析期间接口响应中位数约 0.06 秒\\
 & 逐页流式解析 & 长文档内存占用与页数无关；阅读器按页加载\\
 & Token 感知切分与批量嵌入 & 片段按 token 上限无损拆分；嵌入请求按条数与字节双重限流\\
\midrule
\rtzebra & 检查点与幂等 & 解析、翻译、索引、Agent 规划均可断点续跑；重复提交不重复执行\\
\rtzebra \multirow{-2}{=}{可靠} & 多级降级 & 模型链降级；本地向量 → 云端向量 → 关键词；无密钥模式下解析、阅读、图谱、绘图仍可用\\
\bottomrule
\end{tabularx}
\end{table}
